\textbf{Proof.}
The displayed algebraic coordinates satisfy
\[
\lVert x\rVert_2^2=\lVert y_0\rVert_2^2
=\lVert y_R\rVert_2^2=3,
\qquad
\mathbf1^{\top}x=x^{\top}y_R=0,
\qquad
(\mathbf1^{\top}y_R)^2=\frac{216}{55}.
\tag{1}
\]
Indeed, direct expansion of the two radicals gives \(2t^2+u^2=3\) and \((2t+u)^2=216/55\).  Every displayed coordinate has absolute value below \(2\).  The equality-set formulas in thm:triangle-completion-separation-certificate therefore place \(B_0\) and \(B_R\) in the stated optimizer sets and give
\[
\lambda_{\min}(3^{-1}S_{B_0})
=\lambda_{\min}(3^{-1}S_{B_R})=\frac4{27}>\frac18.
\tag{2}
\]
Thus both bases belong to the certified stable rank-two stratum.  The pilot observations \((0,1)\) in each arm have empirical law \(\operatorname{Bernoulli}(1/2)\), exactly equal to each population marginal.  Their means, variances, quantile-coupling endpoints, and endpoint matrices consequently agree, so both uniform pilot objective errors are zero.  All input coordinates are rational or real algebraic, and thm:encoded-cad-pilot-selector-certificate applies at every positive algebraic tolerance; no agreement with its tie selection is needed for this displayed pair.

For completeness, the finite enumeration underlying the observable variances is the following exact formula.  For \(e\in\{L,U\}\) and \(z\in\{0,1\}^3\), put
\[
Y_i^{(e,z)}(1)=z_i,
\qquad
Y_i^{(U,z)}(0)=z_i,
\qquad
Y_i^{(L,z)}(0)=1-z_i.
\]
For a displayed basis \(B\), let \(T_B(w,z,e)\) be def:observable-coefficient-map after substituting assignment \(w\), array \((e,z)\), the exact triangle \(\Pi,\Omega\), and \(B\).  Then
\[
V_{B,e}
=\frac18\sum_{z\in\{0,1\}^3}
\left[
\sum_{w\in\{0,1\}^3}\nu(w)T_B(w,z,e)^2
-\left\{\sum_{w\in\{0,1\}^3}\nu(w)T_B(w,z,e)\right\}^2
\right].
\tag{3}
\]
Equation (3) is a sum of sixty-four rational or real-algebraic terms.  Exact substitution and common-denominator reduction give the four values in the statement.  Their endpoint differences reduce to
\[
V_{0,L}-V_{0,U}=\frac{137}{1200}>0,
\tag{4}
\]
and
\[
V_{R,L}-V_{R,U}
=\frac{2943215-196776\sqrt{93}}{10750800}>0.
\tag{5}
\]
For (5), \(\sqrt{93}<193/20\) because \(93<37249/400\), and exact cross multiplication makes the numerator positive.  Hence the lower endpoint is worst for both bases.

Subtracting the two lower-endpoint values and, separately, the two exact objective values supplied by thm:triangle-completion-separation-certificate gives the displayed \(G_{\mathrm{obs}}\), \(\delta\), and \(E_{\mathrm{bridge}}\).  The correction is positive because \(\sqrt{93}>241/25\), as \(93>58081/625\), followed by exact cross multiplication.  Finally, rational simplification gives
\[
G_{\mathrm{obs}}-E_{\mathrm{bridge}}
=\frac{4912}{39525}=\delta>0.
\tag{6}
\]
Thus \(G_{\mathrm{obs}}=\delta+E_{\mathrm{bridge}}>E_{\mathrm{bridge}}+0\), where zero is the exact pilot error.  The selected-rank bridge theorem is applicable to either fixed displayed basis by (2); the conclusion makes no statement about other tie representatives or uncontrolled singular ranks. \(\square\)